% ==============================================================================
% Supplementary Methods: Constrained Bivariate Gaussian Mixture Modeling
% Location: 04_Figure_Generation/Supplementary_Figure_06/Supplementary_Methods_GMM_Detailed.tex
% Reference: Bivariate_Gaussian_Mixture_Modeling_Report.tex
% ==============================================================================

\section*{Supplementary Note: Constrained Bivariate Gaussian Mixture Modeling and Model Selection}

\subsection*{1. Biophysical Motivation and Problem Formulation}
Differential Pattern Analysis (DiPa) stratifies genes into functional co-regulation quadrants based on predefined threshold cutoffs ($|\log_2\text{FC}| \ge 0.5$ in mRNA and protein dimensions). A fundamental methodological question is whether the observed coordinate distributions of paired transcriptomic and proteomic shifts arise from a single unimodal bivariate distribution with observational dispersion, or whether the data objectively mandate a finite mixture of distinct generative subpopulations. If a single unimodal distribution were statistically sufficient, the DiPa quadrant taxonomy could be critiqued as a heuristic discretization. Conversely, if multi-component mixture models decisively outperform the single-component null model, the apparent regulatory regimes are established as statistically distinct biological subpopulations.

To resolve this question rigorously and without relying on subjective thresholding, we formulated a parametric framework based on finite bivariate Gaussian Mixture Models (GMM) estimated via the Expectation-Maximization (EM) algorithm, augmented with biophysical structural constraints.

\subsection*{2. Bivariate Mixture Likelihood Formulation}
Let $X_i = (x_i, y_i)^\top \in \mathbb{R}^2$ denote the observed fold changes for gene $i \in \{1, \dots, N\}$ ($N = 943$ core matched pairs analyzed across both BDL and $\text{CCl}_4$ cohorts), where $x_i = \log_2\text{FC}_{\text{RNA}, i}$ and $y_i = \log_2\text{FC}_{\text{Prot}, i}$. Under a finite mixture of $K$ bivariate normal densities, the marginal likelihood of observation $X_i$ given parameter vector $\theta$ is:
\begin{equation}
p(X_i \mid \theta) = \sum_{k=1}^K \pi_k \, \mathcal{N}_2\left(X_i \mid \mu_k, \, \Sigma_k\right),
\label{eq:gmm_marginal}
\end{equation}
subject to the simplex constraints $\pi_k \in [0, 1]$ and $\sum_{k=1}^K \pi_k = 1$, where $\mu_k = (\mu_{xk}, \mu_{yk})^\top \in \mathbb{R}^2$ represents the centroid of component $k$, and $\Sigma_k$ denotes its $2 \times 2$ positive-definite covariance matrix:
\begin{equation}
\Sigma_k = \begin{pmatrix} \sigma_{xk}^2 & \rho_k \sigma_{xk} \sigma_{yk} \\ \rho_k \sigma_{xk} \sigma_{yk} & \sigma_{yk}^2 \end{pmatrix}.
\end{equation}
The bivariate Gaussian density is defined as:
\begin{equation}
\mathcal{N}_2\left(X_i \mid \mu_k, \, \Sigma_k\right) = \frac{1}{2\pi \sqrt{\det \Sigma_k}} \exp\left( -\frac{1}{2} (X_i - \mu_k)^\top \Sigma_k^{-1} (X_i - \mu_k) \right).
\end{equation}
Assuming independent observations across genes, the observed-data log-likelihood is:
\begin{equation}
\ell(\theta) = \sum_{i=1}^N \ln \left( \sum_{k=1}^K \pi_k \, \mathcal{N}_2\left(X_i \mid \mu_k, \, \Sigma_k\right) \right).
\label{eq:incomplete_loglik}
\end{equation}
Because the summation over components appears inside the logarithm, direct analytical maximization of $\ell(\theta)$ is intractable. We therefore introduce latent indicator vectors $Z_i = (Z_{i1}, \dots, Z_{iK})^\top$, where $Z_{ik} \in \{0, 1\}$ and $\sum_{k=1}^K Z_{ik} = 1$, such that $Z_{ik} = 1$ if gene $i$ was generated by component $k$, and $0$ otherwise.

The complete-data log-likelihood is given by:
\begin{equation}
\ell_c(\theta) = \sum_{i=1}^N \sum_{k=1}^K Z_{ik} \left[ \ln \pi_k - \ln(2\pi) - \frac{1}{2}\ln(\det \Sigma_k) - \frac{1}{2} (X_i - \mu_k)^\top \Sigma_k^{-1} (X_i - \mu_k) \right].
\label{eq:complete_loglik}
\end{equation}

\subsection*{3. Expectation-Maximization (EM) Algorithm under Structural Constraints}
Parameters are estimated iteratively by alternating between evaluating the expected value of $\ell_c(\theta)$ with respect to the conditional distribution of $Z_i$ given $X_i$ and the current parameter estimate $\theta^{(t)}$ (E-step), and maximizing this surrogate function over $\theta$ subject to biophysical constraints (M-step).

\subsubsection*{E-Step (Posterior Responsibility Estimation)}
At iteration $t$, the posterior probability (responsibility) $\gamma_{ik}^{(t)} \in [0, 1]$ that component $k$ generated gene $i$ is calculated via Bayes' rule:
\begin{equation}
\gamma_{ik}^{(t)} \equiv \mathbb{E}\left[ Z_{ik} \mid X_i, \, \theta^{(t)} \right] = \frac{\pi_k^{(t)} \, \mathcal{N}_2\left(X_i \mid \mu_k^{(t)}, \, \Sigma_k^{(t)}\right)}{\sum_{j=1}^K \pi_j^{(t)} \, \mathcal{N}_2\left(X_i \mid \mu_j^{(t)}, \, \Sigma_j^{(t)}\right)}.
\label{eq:estep_resp}
\end{equation}
In computation, log-sum-exp stabilization is used to prevent numerical underflow in extreme tail observations:
\begin{equation}
\ln \gamma_{ik}^{(t)} = \ln \pi_k^{(t)} + \ln \mathcal{N}_2\left(X_i \mid \mu_k^{(t)}, \, \Sigma_k^{(t)}\right) - \ln \sum_{j=1}^K \exp\left( \ln \pi_j^{(t)} + \ln \mathcal{N}_2\left(X_i \mid \mu_j^{(t)}, \, \Sigma_j^{(t)}\right) \right).
\end{equation}

\subsubsection*{M-Step (Constrained Maximization)}
Defining the effective sample size belonging to component $k$ as $N_k^{(t)} = \sum_{i=1}^N \gamma_{ik}^{(t)}$, the unconstrained maximum-likelihood updates for mixing proportions and centroid vectors are:
\begin{equation}
\pi_k^{(t+1)} = \frac{N_k^{(t)}}{N}, \qquad \mu_k^{(t+1)} = \frac{1}{N_k^{(t)}} \sum_{i=1}^N \gamma_{ik}^{(t)} X_i.
\end{equation}
The weighted sample dispersion moments are evaluated as raw empirical quantities:
\begin{align}
\sigma_{xk}^{2, \text{raw}} &= \frac{1}{N_k^{(t)}} \sum_{i=1}^N \gamma_{ik}^{(t)} \left(x_i - \mu_{xk}^{(t+1)}\right)^2, \\
\sigma_{yk}^{2, \text{raw}} &= \frac{1}{N_k^{(t)}} \sum_{i=1}^N \gamma_{ik}^{(t)} \left(y_i - \mu_{yk}^{(t+1)}\right)^2, \\
\text{Cov}_{xy, k}^{\text{raw}} &= \frac{1}{N_k^{(t)}} \sum_{i=1}^N \gamma_{ik}^{(t)} \left(x_i - \mu_{xk}^{(t+1)}\right)\left(y_i - \mu_{yk}^{(t+1)}\right), \qquad \rho_k^{\text{raw}} = \frac{\text{Cov}_{xy, k}^{\text{raw}}}{\sigma_{xk}^{\text{raw}} \sigma_{yk}^{\text{raw}}}.
\end{align}

\subsubsection*{Biophysical Structural Constraints}
To align components directly with hypothesized regulatory regimes and prevent auxiliary components from fitting local unstructured point noise, structural constraints are enforced:
\begin{enumerate}
    \item \textbf{Component 1: Stoichiometric Co-Regulation (Diagonal Orientation).}  
    Captures canonical proportional transcription--translation coupling. The Pearson correlation is restricted to a positive interval reflecting strong, biologically plausible coupling:
    \begin{equation}
    \rho_1^{(t+1)} = \min\left(0.90, \, \max\left(0.70, \, \rho_1^{\text{raw}}\right)\right).
    \end{equation}
    The covariance matrix is reconstituted as:
    \begin{equation}
    \Sigma_1^{(t+1)} = \begin{pmatrix} \sigma_{x1}^2 & \rho_1^{(t+1)} \sigma_{x1} \sigma_{y1} \\ \rho_1^{(t+1)} \sigma_{x1} \sigma_{y1} & \sigma_{y1}^2 \end{pmatrix}.
    \end{equation}

    \item \textbf{Component 2: Protein-Dominant Decoupling (Vertical Orientation).}  
    Captures post-transcriptional divergence where protein abundance changes markedly with minimal transcriptomic shift. We enforce strict coordinate orthogonality ($\rho_2 \equiv 0$) and vertical elongation controlled by minimum variance ratio factor $c$:
    \begin{equation}
    \rho_2^{(t+1)} = 0, \quad \sigma_{x2}^2 = \sigma_{x2}^{2, \text{raw}}, \quad \sigma_{y2}^2 = \max\left(\sigma_{y2}^{2, \text{raw}}, \, c \cdot \sigma_{x2}^2\right) \implies \Sigma_2^{(t+1)} = \begin{pmatrix} \sigma_{x2}^2 & 0 \\ 0 & \sigma_{y2}^2 \end{pmatrix}.
    \end{equation}

    \item \textbf{Component 3: Post-Transcriptional Buffering (Horizontal Orientation, for $K=3$).}  
    Captures translational/transcriptomic buffering where large mRNA shifts fail to propagate to the protein level. We enforce strict orthogonality ($\rho_3 \equiv 0$) and horizontal elongation:
    \begin{equation}
    \rho_3^{(t+1)} = 0, \quad \sigma_{y3}^2 = \sigma_{y3}^{2, \text{raw}}, \quad \sigma_{x3}^2 = \max\left(\sigma_{x3}^{2, \text{raw}}, \, c \cdot \sigma_{y3}^2\right) \implies \Sigma_3^{(t+1)} = \begin{pmatrix} \sigma_{x3}^2 & 0 \\ 0 & \sigma_{y3}^2 \end{pmatrix}.
    \end{equation}
\end{enumerate}
Iterations proceed until convergence: $|\ell(\theta^{(t+1)}) - \ell(\theta^{(t)})| < 10^{-6}$ (typically achieved within 50--200 iterations).

\subsection*{4. Exact Enumeration of Free Parameters ($P$)}
Information-theoretic model selection requires exact counting of freely estimated parameters:
\begin{itemize}
    \item \textbf{Model $K = 1$ (Null Single Gaussian):} 2 centroid coordinates ($\mu_x, \mu_y$) + 2 variances ($\sigma_x^2, \sigma_y^2$) + 1 correlation ($\rho$) $\implies P = 5$.
    \item \textbf{Model $K = 2$ (Constrained Mixture):} 1 independent mixing weight ($\pi_1$, with $\pi_2 = 1 - \pi_1$) + 4 centroid coordinates ($\mu_1, \mu_2$) + 3 dispersion parameters for Comp~1 ($\sigma_{x1}^2, \sigma_{y1}^2, \rho_1$) + 2 independent variances for Comp~2 ($\sigma_{x2}^2, \sigma_{y2}^2$, since $\rho_2 \equiv 0$) $\implies P = 1 + 4 + 3 + 2 = 10$.
    \item \textbf{Model $K = 3$ (Constrained Mixture):} 2 independent mixing weights ($\pi_1, \pi_2$) + 6 centroid coordinates ($\mu_1, \mu_2, \mu_3$) + 3 dispersion parameters for Comp~1 + 2 variances for Comp~2 ($\rho_2 \equiv 0$) + 2 variances for Comp~3 ($\rho_3 \equiv 0$) $\implies P = 2 + 6 + 3 + 2 + 2 = 15$.
\end{itemize}

\subsection*{5. Model Selection via Bayesian Information Criterion (BIC)}
Model selection across component numbers $K \in \{1, 2, 3\}$ and elongation factors $c \in \{1.3, 1.5, 2.0, 2.5, 4.0\}$ was conducted using the Bayesian Information Criterion:
\begin{equation}
\text{BIC} = -2\ln \hat{L} + P \ln(N),
\label{eq:bic}
\end{equation}
where $\hat{L} = \exp(\ell(\hat{\theta}))$ is the maximized likelihood and $N = 943$. The penalty term for each additional parameter is $\ln(943) \approx 6.849$, imposing a stringent barrier against over-parameterization. Improvements are measured via $\Delta\text{BIC} = \text{BIC}_{\text{Model}} - \text{BIC}_{K=1}$. Following the Kass and Raftery scale of evidence, a difference $\Delta\text{BIC} < -10$ indicates decisive statistical evidence against the simpler model.

\subsubsection*{Theoretical Justification for BIC over Likelihood-Ratio Tests (LRT)}
Classical Likelihood-Ratio Tests ($\Lambda = 2(\ell_1 - \ell_0)$) assume that under $H_0$, the test statistic asymptotically follows a $\chi^2$ distribution with degrees of freedom equal to the parameter difference (Wilks' theorem). However, in finite mixture models, testing $K$ against $K+1$ components corresponds to setting the null hypothesis $\pi_{K+1} = 0$. This hypothesis lies strictly on the \textit{boundary} of the probability simplex ($\pi \in [0, 1]^K$), rather than in its interior. On the boundary, the score function vanishes, the Fisher information matrix is singular, and the standard regularity conditions of Wilks' theorem are violated. Consequently, asymptotic $\chi^2$ distributions fail to hold, and LRT p-values are invalid. The BIC is asymptotically consistent and robust to boundary violations, making it the gold standard for mixture complexity selection.

\subsection*{6. Sensitivity Analysis across Elongation Constraints}
To verify that the chosen component complexity was not an artifact of a specific geometric threshold, the elongation constraint factor $c$ was evaluated across five values: $c \in \{1.3, 1.5, 2.0, 2.5, 4.0\}$, corresponding to geometric semi-axis aspect ratios $\sqrt{c} \in \{1.14, 1.22, 1.41, 1.58, 2.00\}$.

As documented in \textbf{Supplementary Table~5}:
\begin{enumerate}
    \item \textbf{Decisive rejection of single Gaussian null:} In both cohorts and across all values of $c$, the single Gaussian model ($K=1$) was conclusively rejected ($\Delta\text{BIC} \le -259.7$ for BDL; $\Delta\text{BIC} \le -106.5$ for $\text{CCl}_4$).
    \item \textbf{Invariance of optimal complexity in BDL ($K = 2$):} In BDL, $K = 2$ consistently achieved lower BIC than $K = 3$ across all values of $c$ (e.g., at $c = 1.3$: $\text{BIC}_{K=2} = 3390.5$ vs $\text{BIC}_{K=3} = 3398.5$, $\Delta\text{BIC} = +8.0$), proving that a third buffering component is statistically unneeded in cholestasis.
    \item \textbf{Invariance of optimal complexity in $\text{CCl}_4$ ($K = 3$):} In $\text{CCl}_4$, $K = 3$ decisively outperformed $K = 2$ across all tested constraints (e.g., at $c = 1.3$: $\text{BIC}_{K=3} = 3490.6$ vs $\text{BIC}_{K=2} = 3606.4$, $\Delta\text{BIC} = -115.8$), demonstrating that transcriptomic buffering is an indispensable biological regime under toxic necrosis.
\end{enumerate}

\subsection*{7. Correspondence between GMM Components and DiPa Quadrants}
The fitted Gaussian components demonstrate direct biological and geometrical mapping to the DiPa taxonomy:
\begin{itemize}
    \item \textbf{Component 1 (Diagonal, Blue):} Corresponds to DiPa Quadrants~1 (concordant induction), 2 (concordant repression), and Quadrant~8 (constitutive baseline co-expression).
    \item \textbf{Component 2 (Vertical, Purple):} Corresponds to DiPa Quadrants~3 (isolated protein accumulation) and 4 (isolated protein depletion), representing protein-dominant proteostatic decoupling.
    \item \textbf{Component 3 (Horizontal, Orange):} Corresponds to DiPa Quadrants~5 (isolated mRNA induction) and 6 (isolated mRNA repression), representing transcriptomic buffering.
\end{itemize}
This concordance confirms that the DiPa functional regimes reflect bona fide statistical subpopulations governed by distinct post-transcriptional regulatory mechanisms.
